\DocumentMetadata{
  testphase = {phase-II, math},
  activate = tagging,
  pdfstandard = A-2b,
  pdfversion = 1.7
}
\documentclass{article}
\usepackage[a4paper,margin=2cm]{geometry}
\usepackage{hyperref}
\title{atscv Manual}
\date{2026-04-29}
\begin{document}
\maketitle
\section{Introduction}
atscv is a structured CV compiler for ATS parsing and human readability.
\section{ATS principles}
Linear reading order, semantic headings, plain bullets, and predictable tokenization are mandatory.
\section{Why LaTeX fails in ATS (and how atscv fixes it)}
Complex layouts (tables, columns, sidebars) scramble extracted text. atscv avoids those structures and enforces one flow.
\section{Design modes explained}
ats, modern, executive, consulting, and compact provide controlled styling while preserving ATS-safe structure.
\section{How to write a good CV (STAR + impact)}
Use Situation, Task, Action, Result and quantify outcomes by KPI, time, or cost.
\section{Human scan principles (30-second rule)}
Top-load impact, maintain hierarchy, and keep bullets concise and outcome-first.
\section{Multilingual usage}
LuaLaTeX + Unicode enables multilingual rendering without disrupting extraction order.
\section{Example walkthrough}
Compile examples/fabio-bigtech-modern.tex, run pdftotext, and verify ordering.
\section{ATS validation workflow}
Build with LuaLaTeX, check log warnings, run pdftotext, and compare extracted text to source intent.
\end{document}
